\documentclass[11pt]{article}

\usepackage[margin=0.9in]{geometry}
\usepackage{amsmath}
\usepackage{amssymb}
\usepackage{booktabs}
\usepackage{graphicx}
\usepackage{float}
\usepackage[hidelinks]{hyperref}
\usepackage{array}
\usepackage[table]{xcolor}
\usepackage{parskip}
\usepackage{mdframed}

\graphicspath{{../figures/}}
\sloppy
\setcounter{tocdepth}{2}

\newcommand{\mainfigure}[3]{%
\begin{figure}[H]
    \centering
    \includegraphics[width=\textwidth]{#1}
    \caption{#2}
    \label{#3}
\end{figure}
}

\newmdenv[backgroundcolor=yellow!12, linecolor=gray!60, linewidth=0.6pt,
          innertopmargin=6pt, innerbottommargin=6pt]{checkenv}

\title{Interpreting FCFS Appointment Simulation Results\\
Under the Current Formulation\\[0.4em]
\large Metric behavior, access losses, and mechanisms}
\author{CUIMC Appointment Simulation}
\date{\today}

\begin{document}
\maketitle

\begin{abstract}
\noindent
This note interprets simulation outputs from a discrete-event FCFS appointment
model. The goal is to understand what the outputs show and whether they are
mechanically plausible under the current formulation---not to make empirical
claims about real clinics.

The central finding is that utilization and served rate separate sharply.
At baseline ($\lambda = 100$ arrivals per day, 32 slots per day, 14-day
horizon), average utilization is 0.839 while overall served rate is 0.269.
This separation is mechanically plausible because the two metrics have
different denominators: utilization is normalized by slots, served rate by
arrivals. High slot fill does not imply that most patients receive care.

Arrival outcome decomposition at baseline shows that balking (35.6\% of
arrivals) and cancellation (32.3\%) are the dominant loss channels, not
horizon saturation: the no-offer rate is~0 at baseline demand. The high
cancellation fraction is mechanically consistent with the long average
booking delay (8.35~days) combined with a 10\%~daily cancellation
probability.

Mean offered wait is preferable to mean accepted wait as a delay measure,
because offered delay is recorded before the balking decision and therefore
includes patients who ultimately reject the offer. Offered wait still excludes
no-offer patients, so it must be read together with the no-offer share---which
is zero at baseline but grows above $\lambda \approx 110$.
\end{abstract}

\tableofcontents
\newpage

%% ============================================================
\section{Simulation Mechanics Relevant for Interpretation}
\label{sec:mechanics}
%% ============================================================

The following mechanics directly affect how outputs should be read. The full
formulation is in the simulation documentation.

\paragraph{Finite booking horizon and FCFS pooling.}
Each day, arrivals for both classes are pooled and randomly permuted into one
FCFS sequence. Each patient is offered the earliest available slot within the
$H$-day horizon (including same-day, $r=0$). Because FCFS operates on a
permuted joint sequence, class labels are irrelevant to slot assignment:
neither class is prioritized. Class disparities in outcomes arise only from
differences in arrival rates or behavioral parameters.

\paragraph{Offered delay $\tau$ recorded before balking.}
When a valid slot exists, the offered delay $\tau$ is recorded before the
patient decides whether to accept. Mean offered wait therefore includes
patients who ultimately balk. Mean accepted wait excludes them, which biases
it downward when high-$\tau$ patients disproportionately balk.

\paragraph{No-show based on original $\tau$.}
No-show probability is a function of the original offered delay $\tau$, not of
the remaining residual day $r$ at service time. A patient who accepted a slot
9~days out carries their high-delay no-show probability regardless of when they
end up being seen.

\paragraph{Future cancellations only; no rebooking of no-show slots.}
Cancellations apply only to future appointments ($r \geq 1$). Same-day
cancellations are not modeled. Slots vacated by no-shows are not
rebooked---once a patient no-shows, that slot is lost.

\paragraph{No-offer rule.}
If the horizon is full when a patient reaches the front of the FCFS queue,
that patient and all remaining arrivals for that day are counted as
\texttt{no\_offer}. This means a single saturation event affects the entire
remainder of the day's queue, which could introduce daily-order artifacts.

\paragraph{Utilization denominator.}
Utilization counts completed visits per available slot, not scheduled
appointments per slot. No-shows, cancellations, and unbooked slots all reduce
utilization. A high-utilization day is one where most slots became completed
visits.

%% ============================================================
\section{Baseline Diagnosis}
\label{sec:baseline}
%% ============================================================

\subsection{Parameter Setup}

The baseline (\texttt{configs/baseline.yaml}) uses 32~slots per day,
a 14-day booking horizon, two symmetric classes with $\lambda_1 = \lambda_2
= 50$~arrivals per day, cancellation probability~0.10, balking threshold
9~days with high-delay balking probability~0.50, and no-show threshold 6~days
with high-delay no-show probability~0.30. Low-delay balking and no-show
probabilities are~0. Measurement window is 365~days after a 30-day burn-in.

The baseline parameters are intentionally stress-testing the simulation:
100~arrivals per day against 32~slots per day is a 3:1 demand-to-capacity
ratio, well above any realistic clinic load. The purpose is to probe how the
metrics behave across a wide range, not to model a plausible clinic.

\subsection{Aggregate Metrics at Baseline}

\begin{table}[H]
\centering
\renewcommand{\arraystretch}{1.15}
\begin{tabular}{lr}
\toprule
Metric & Value \\
\midrule
Average utilization & 0.839 \\
Overall served rate & 0.269 \\
Mean offered wait (days) & 9.30 \\
Mean accepted wait (days) & 8.35 \\
Overall balking rate & 0.356 \\
Class served-rate gap ($\Delta_{\text{access}}$) & 0.001 \\
\bottomrule
\end{tabular}
\caption{Baseline aggregate metrics. The near-zero class gap is expected:
both classes have identical behavioral parameters and symmetric arrival rates,
so class labels are exchangeable.}
\label{tab:baseline}
\end{table}

\subsection{The Utilization--Access Gap}

Utilization~(0.839) and served rate~(0.269) differ by a factor of roughly
3.1. This is mechanically expected. Utilization is
$\text{served slots} / \text{slots per day}$, normalized by supply.
Served rate is $\text{served patients} / \text{total arrivals}$, normalized
by demand. When demand far exceeds supply, served rate is bounded by the
supply-to-demand ratio; utilization is not.

At baseline, $100/32 = 3.125$~arrivals per slot. Even if every slot were
filled, served rate could not exceed $1/3.125 = 0.32$. The observed
served rate of 0.269 is slightly below this ceiling because some filled slots
are lost to no-shows.

The simulation implies, under these assumptions, that a facility-level
utilization measure would give a misleading picture of patient access at
this demand-to-capacity ratio.

%% ============================================================
\section{Arrival Outcome Decomposition}
\label{sec:decomp}
%% ============================================================

\subsection{Loss Channels at Baseline}

Every arrival ends up in exactly one outcome state. At baseline:

\begin{table}[H]
\centering
\renewcommand{\arraystretch}{1.12}
\begin{tabular}{lrp{0.52\textwidth}}
\toprule
Outcome & Share of arrivals & Interpretation \\
\midrule
Served & 0.269 & Completed visit \\
Balked & 0.356 & Received an offer, rejected the delay \\
Canceled & 0.323 & Booked, but canceled before service day \\
No-show & 0.052 & Kept booking, did not attend \\
No-offer & 0.000 & Horizon full; no slot offered \\
Unresolved & $<$0.001 & Booked but not resolved before simulation end \\
\midrule
Total & 1.000 & \\
\bottomrule
\end{tabular}
\caption{Arrival outcome decomposition at baseline ($\lambda = 100$).
No-offer is zero: the horizon does not completely fill at baseline demand
because cancellations continuously free future capacity.}
\label{tab:decomp}
\end{table}

Two findings stand out. First, the no-offer rate is zero even at 100
arrivals/day. The horizon does not saturate because cancellations
(10\%/day for each future appointment) continuously free slots. However,
cancellation frees far-future capacity, which is then offered to new
arrivals at long delays---reinforcing the same long-delay dynamic that
drives balking.

Second, the 32.3\% cancellation rate (as a share of arrivals) is not
self-evidently wrong. Of the 64.4\% of arrivals that book an appointment,
roughly half ultimately cancel. With an average booking delay of 8.35~days
and a 10\%~daily cancellation probability, the implied cumulative
cancellation probability is $1 - 0.9^{8.35} \approx 0.59$. The observed
50\% cancellation rate among booked patients is in the same range, though
somewhat lower---possibly because some booked patients are scheduled for
soon-enough slots that they face fewer cancellation check days.

\begin{checkenv}
\textbf{Check.} The high cancellation fraction warrants a manual trace.
The cumulative cancellation argument assumes each booked patient faces
a fresh 10\%~draw every day until their appointment. If the engine applies
cancellation only once per booking rather than once per residual day, the
observed rate would be much lower. The current documentation states
cancellations apply to ``future appointments'' each day, which implies
once per residual day, but this should be confirmed in
\texttt{simulation/engine.py}.
\end{checkenv}

\subsection{How Loss Channels Shift with Demand}

Figure~\ref{fig:decomp} shows the arrival outcome decomposition as total
daily arrivals vary from 30 to 170. Several patterns are mechanically
consistent with the formulation.

At low demand ($\lambda \lesssim 50$), virtually all arrivals are served.
The horizon is sparse, offered delays are short, balking is near zero,
and cancellation---although it still happens---mostly frees slots that
are quickly refilled. As demand rises above the 32-slot daily capacity,
the served rate falls. The no-offer rate remains zero up to $\lambda
\approx 110$, at which point the horizon begins to saturate. Above
$\lambda = 110$, no-offer becomes an additional and mechanically distinct
loss channel.

The transition from a cancellation- and balking-dominated regime (below
$\lambda \approx 110$) to a horizon-saturation regime (above) is a
structural feature of the model. It suggests that interpreting no-offer
share alongside served rate and offered wait is necessary to understand
which failure mode dominates in a given parameter setting.

\mainfigure{research_note_main_panel.png}{%
\textit{Left:} Arrival outcome decomposition as total daily arrivals
vary. At baseline ($\lambda = 100$, dashed line), balking and
cancellation are the dominant loss channels; no-offer is
zero. No-offer appears only above $\lambda \approx 110$.
\textit{Center:} Utilization--access frontier. Utilization is nearly
flat at~0.84 across all demand levels, while served rate ranges from
near~1 to below~0.16. These two metrics are essentially independent.
\textit{Right:} Offered vs.\ accepted wait. Points below the~45$^\circ$
line indicate that accepted wait understates the delay patients face.
The gap widens with balking rate (red intensity).}{fig:main-panel}

%% ============================================================
\section{Metric Hierarchy Under This Formulation}
\label{sec:metrics}
%% ============================================================

\begin{table}[H]
\centering
\renewcommand{\arraystretch}{1.2}
\begin{tabular}{p{0.20\textwidth}p{0.22\textwidth}p{0.24\textwidth}p{0.24\textwidth}}
\toprule
Metric & What it measures & Mechanical interpretation & Main caveat \\
\midrule
Overall served rate & Arrivals $\to$ completed visits & Primary access measure; bounded by supply/demand ratio & Does not decompose loss channels \\
Average utilization & Completed visits per slot & Capacity-use measure; independent of demand level when system is saturated & High value does not imply good access \\
Mean offered wait & Delay among patients who received any offer & Less selective than accepted wait; includes patients who then balk & Excludes no-offer patients entirely \\
Mean accepted wait & Delay among patients who accepted bookings & Conditional; pulled down when high-delay patients balk & Can fall as access worsens \\
Balking rate & Share of offered patients who reject & Congestion diagnostic & Not a success metric \\
Class served-rate gap & Class 1 served rate minus Class 2 & Meaningful only when classes differ in behavior or arrivals & Near-zero in symmetric baseline \\
\bottomrule
\end{tabular}
\caption{Metric interpretation under the current FCFS formulation.}
\label{tab:metrics}
\end{table}

\subsection{Utilization and Access Are Structurally Decoupled}

The center panel of Figure~\ref{fig:main-panel} shows that utilization is
essentially constant at~0.839--0.841 across all demand levels from
$\lambda = 60$ to~170, while served rate falls monotonically from 0.45 to
0.16. This is not an accidental numerical coincidence---it is a direct
consequence of how the model fills slots. As demand rises, more patients
compete for each slot, more balk at long delays, and the survivors who book
tend to be more likely to eventually cancel; the net effect keeps the slot-fill
rate near the no-show--adjusted ceiling. Served rate, normalized by the much
larger arrival pool, falls continuously.

\subsection{Offered Wait vs.\ Accepted Wait}

The right panel of Figure~\ref{fig:main-panel} shows that accepted wait
(y-axis) is consistently below offered wait (x-axis), placing all points
below the 45$^\circ$ line. The gap widens as balking rate rises. This is
mechanically consistent with the formulation: patients who reject offers
have above-average $\tau$, so removing them from the accepted sample
pulls the average down.

The simulation implies, under these assumptions, that using accepted wait as
a patient-facing delay measure would understate congestion in high-balking
regimes. Offered wait is the better denominator for any congestion
diagnostic---but it still excludes no-offer patients, who receive no delay
information at all.

\subsection{No-Offer as a Distinct Access Failure}

Mean offered wait is computed among patients who received a slot offer.
When the horizon saturates ($\lambda \gtrsim 110$ in this formulation),
a growing share of arrivals receive no offer and are entirely absent from
the offered-wait calculation. At $\lambda = 170$, the no-offer share is
34.6\% of arrivals. Reading offered wait alone in this regime misses the
worst-off patients.

A useful diagnostic is to report no-offer share alongside offered wait.
Together they describe both the level and the completeness of delay
information.

%% ============================================================
\section{Sensitivity Screen: Main Patterns}
\label{sec:sensitivity}
%% ============================================================

The one-at-a-time sensitivity runs and the regression screen are described
in the supplementary appendix. The patterns most relevant for
interpretation are summarized here.

\paragraph{Demand pressure.}
Total arrival rate has the largest absolute standardized regression
coefficient for overall served rate ($\beta = -0.774$) and is the
dominant driver of offered wait ($\beta = +0.576$). This is mechanically
consistent: more arrivals compete for a fixed number of slots, pushing
offers farther out and reducing access.

\paragraph{No-show and cancellation.}
Average no-show threshold has the largest coefficient for utilization
($\beta = +0.512$): a higher threshold delays the onset of high no-show
probability, keeping more booked slots as completed visits. Cancellation
has a positive coefficient for served rate ($\beta = +0.117$) and a
negative one for offered wait ($\beta = -0.441$). The positive sign on
served rate is counterintuitive at first: the mechanism is that
cancellations free future capacity, which can absorb more patients and
slightly improve the served fraction---a rebooking-chain effect under FCFS.

\begin{checkenv}
\textbf{Check.} The positive cancellation coefficient on served rate
should be verified. It implies that, in the simulation, more cancellation
slightly helps overall access. This depends critically on whether
canceled slots are immediately rebooked by new arrivals (which the FCFS
formulation allows) and whether the net effect on served rate is
positive after accounting for the original canceled patient. This should
be checked against the arrival outcome decomposition at varying
cancellation probabilities.
\end{checkenv}

\paragraph{Balking.}
Higher balking step ($\beta = -0.206$ on offered wait) reduces offered
wait through a selection effect: patients with long $\tau$ leave the
accepted sample and eventually the offered sample shrinks toward shorter
delays. This is not a congestion improvement. The deep-dive analysis
(Section~\ref{sec:balk-deepdive}) confirms that accepted wait can fall
while served rate also falls.

\paragraph{Class gap drivers.}
In the randomized screen, the largest class-gap drivers are cancellation
probability gap ($|\beta| = 0.452$), balking threshold gap ($|\beta| =
0.429$), and balking step gap ($|\beta| = 0.349$). These findings are
plausible: in a pooled FCFS system, class disparities arise because one
class accepts more long-delay offers or is more likely to cancel after
booking.

\mainfigure{regression_significant_standardized_coefficients.png}{%
Significant standardized regression coefficients from 240 randomized
parameter settings. Features and targets are standardized. The
dominance of total arrival rate for served rate and offered wait is
consistent with the structural decoupling of utilization and access.
Error bars are 95\% HC3 robust confidence intervals.}{fig:regression}

%% ============================================================
\section{Balking Deep Dive: Selection vs.\ Access}
\label{sec:balk-deepdive}
%% ============================================================

This section isolates Class~1 balking while Class~2 stays at baseline, to
test whether lower accepted wait can accompany lower access.

\mainfigure{../../../class1_balking/figures/percent_serviced_by_class.png}{%
As Class~1 high-balking probability increases from 0 to~1, Class~1
served rate falls while Class~2 served rate rises. The congestion
relief that Class~2 gains is mechanically consistent with fewer
Class~1 patients booking long-delay slots.}{fig:balk-served}

\mainfigure{../../../class1_balking/figures/mean_accepted_delay_by_class.png}{%
Class~1 accepted wait falls as Class~1 high-balking probability
increases. This is a selection artifact: long-delay Class~1 patients
leave the accepted pool. Accepted wait should not be read as a
delay-improvement signal in this context.}{fig:balk-accepted}

The two figures together demonstrate a pattern the formulation produces
cleanly: accepted wait and served rate move in opposite directions
when balking changes. Any policy that raises balking step would show
apparently improved wait times while actually reducing the share of
Class~1 patients who complete service.

%% ============================================================
\section{Class-Gap Interpretation Under FCFS}
\label{sec:classgap}
%% ============================================================

Arrivals from both classes are pooled and randomly permuted each day before
FCFS assignment. The engine does not distinguish between classes when
offering slots. Therefore:

\begin{itemize}
\item When class parameters are symmetric, the expected class served-rate
gap is zero. Observed deviations are simulation noise.
\item Meaningful class gaps arise only from asymmetric behavioral parameters
(different balking threshold, balking step, cancellation probability, or
no-show behavior) or from different arrival rates. These create systematic
differences in how often each class successfully books, cancels, and
attends.
\item A positive class~1 balking threshold relative to class~2 means
class~1 tolerates longer delays before rejecting offers. Under FCFS, a
more tolerant class captures more of the later-horizon slots and therefore
achieves a higher served rate.
\end{itemize}

\mainfigure{metric_class_gap_drivers.png}{%
When Class~1 behavioral parameters are varied while Class~2 stays at
baseline, a gap opens. Higher Class~1 cancellation, higher balking step,
or higher no-show step widen the gap against Class~1. Higher Class~1
balking threshold helps Class~1 by allowing it to accept longer-delay
slots. With symmetric parameters, both class lines coincide.
}{fig:class-gap}

The regression screen confirms this ranking: among class-gap drivers,
cancellation probability gap dominates ($|\beta| = 0.452$), followed by
balking threshold gap ($|\beta| = 0.429$) and balking step gap
($|\beta| = 0.349$). Class~1 arrival share enters significantly
($\beta = -0.120$): a larger Class~1 share tends to mildly reduce the
Class~1 served-rate advantage, consistent with more Class~1 arrivals
competing for the same slots under FCFS.

%% ============================================================
\section{What Seems Mechanically Robust vs.\ What Needs Validation}
\label{sec:checks}
%% ============================================================

\subsection{Mechanically Robust Under the Current Formulation}

\begin{itemize}
\item \textbf{Utilization and served rate have different denominators.}
Their decoupling is a direct consequence of the metric definitions. It
is not a fragile numerical result.

\item \textbf{High utilization can coexist with low served rate.}
When demand exceeds supply, served rate is bounded by the supply-to-demand
ratio. This bound holds regardless of behavioral parameters.

\item \textbf{Offered wait is less selective than accepted wait.}
Offered delay is recorded before balking, so offered wait includes patients
who ultimately leave. Accepted wait excludes them. The direction of the
difference is structural.

\item \textbf{No-offer is a separate access failure not captured by
offered wait.} No-offer patients receive no delay information and are
excluded from all delay metrics. Above $\lambda \approx 110$, they become
a substantial fraction of arrivals.

\item \textbf{Class labels are exchangeable under symmetric parameters.}
FCFS pooling and random permutation mean that class identity has no
systematic effect when classes are identical. The near-zero baseline class
gap (0.001) is consistent with this.

\item \textbf{The cancellation--delay feedback is directionally correct.}
Longer average booking delay creates more cancellation-check opportunities,
which raises cumulative cancellation probability. The high observed
cancellation fraction at long delays is mechanically consistent with this.
\end{itemize}

\subsection{Needs Validation or Further Code Checks}

\begin{checkenv}
\begin{enumerate}
\item \textbf{Confirm cancellation applies once per residual day.}
The documentation says cancellation is checked each day for all future
appointments. If the engine instead applies cancellation once at booking,
the high cancellation fraction is unexplained and a code audit is needed.
Check \texttt{simulation/engine.py} cancellation loop.

\item \textbf{Clarify no-offer propagation within a day.}
The documentation states that when the horizon is full for one patient,
all remaining arrivals that day are counted as \texttt{no\_offer}. This
means a single saturation event affects the full day remainder, which
could create day-order artifacts. Verify whether the no-offer share is
sensitive to the random permutation seed.

\item \textbf{Verify positive cancellation coefficient on served rate.}
The regression screen finds that average cancellation probability has a
positive standardized coefficient on overall served rate ($\beta = +0.117$).
The proposed mechanism (freed slots rebooked by new arrivals) should be
verified by inspecting the arrival outcome decomposition at varying
cancellation probabilities.

\item \textbf{Assess whether no-show slots should remain unbooked.}
No-show slots are not rebooked in the current formulation. In many
real clinics, same-day no-shows are filled with walk-ins or advance
overbooking. If this assumption is changed, utilization would rise but
served rate would also change---the magnitude depends on the no-show
fraction and same-day availability.

\item \textbf{Assess whether same-day cancellations should be excluded.}
Currently, no same-day cancellations are allowed. This keeps service-day
slots deterministic once the service day is reached. Relaxing this would
add a slot-recapture mechanism that could affect both utilization and
served rate.

\item \textbf{Confirm stability of class-gap estimates.}
The class-gap metric has the lowest out-of-sample $R^2$ (0.708) in the
regression screen and is near zero (0.001) at symmetric baseline. With
only 2~seeds per setting in the regression screen, sampling variance in
the gap could be material. At least 10--20 seeds per grid point are
advisable before drawing strong conclusions about class equity.

\item \textbf{Evaluate realism of baseline demand and capacity levels.}
100 arrivals per day against 32 slots per day is an intentional
stress-test, not a calibrated clinic. Before drawing any operational
conclusions, behavioral parameters (balking thresholds, no-show
probabilities, cancellation rates) should be estimated from clinic data
and the demand-to-capacity ratio brought into a plausible range.
\end{enumerate}
\end{checkenv}

%% ============================================================
\section{Summary of Key Patterns}
\label{sec:summary}
%% ============================================================

\begin{enumerate}
\item Under the current formulation, utilization and served rate are
structurally decoupled. At baseline, this produces a 3:1 discrepancy
(0.839 vs.\ 0.269). Using utilization as a proxy for access would be
misleading at this demand level.

\item The dominant loss channels at baseline are balking (35.6\%) and
cancellation (32.3\%). No-offer is zero. Above $\lambda \approx 110$,
horizon saturation becomes an additional channel.

\item The high cancellation fraction is consistent with the long average
booking delay and the daily-cancellation rule. Whether this is realistic
depends on how real patients behave when booked many days out.

\item Offered wait is preferable to accepted wait as a delay measure,
because it includes patients who rejected the offer. But it still excludes
no-offer patients, who are the most severely access-constrained.

\item Class disparities in the simulation arise from asymmetric behavioral
parameters, not from FCFS favoritism. Symmetric classes produce
near-zero gaps.

\item The regression screen and one-at-a-time sensitivity runs are
internally consistent: demand pressure drives served rate and offered
wait; no-show and cancellation behavior drive utilization; behavioral
parameter gaps drive class equity.

\item None of these findings should be interpreted as causal estimates
about real appointment systems. They describe how this particular model
formulation responds to parameter variation.
\end{enumerate}

%% ============================================================
\appendix
\section{Supplementary Sensitivity Plots}
\label{app:sens}
%% ============================================================

\subsection{Driver Plots}

\mainfigure{metric_utilization_drivers.png}{%
Utilization across Class~1 arrival and behavioral parameter ranges.
Utilization stays near 0.839 across most of the demand range, while
class-level served rates diverge.}{fig:util-drivers}

\mainfigure{metric_access_drivers.png}{%
Served rate falls as arrival rates rise. The Class~1 and Class~2 lines
separate when Class~1 parameters diverge from symmetric baseline.
}{fig:access-drivers}

\mainfigure{metric_wait_drivers.png}{%
Offered and accepted wait across Class~1 parameter ranges. Accepted
wait can fall while offered wait stays elevated, illustrating selection
bias in the accepted-wait metric.}{fig:wait-drivers}

\mainfigure{metric_balking_rate_drivers.png}{%
Balking rate is controlled primarily by balking step and threshold.
Rising balking rate indicates congestion but not improved service.
}{fig:balk-drivers}

\subsection{No-Show and Cancellation Heatmaps}

\mainfigure{no_show_step_interaction_heatmap.png}{%
No-show step heatmap. Utilization falls when both classes carry high
no-show risk. Class gap moves against the higher no-show class.
}{fig:noshow-step}

\mainfigure{no_show_threshold_interaction_heatmap.png}{%
No-show threshold heatmap. Higher threshold delays high no-show onset,
improving utilization and access for the affected class.
}{fig:noshow-threshold}

\mainfigure{cancellation_probability_interaction_heatmap.png}{%
Cancellation heatmap. Higher cancellation for one class reduces that
class's access; aggregate effects are more complex (see
Section~\ref{sec:sensitivity}).}{fig:cancel}

\subsection{Balking Heatmaps}

\mainfigure{balking_step_interaction_heatmap.png}{%
Balking step heatmap. Served rate falls and class gap widens as one
class's balking step rises above the other.}{fig:balk-step}

\mainfigure{balking_threshold_interaction_heatmap.png}{%
Balking threshold heatmap. The more tolerant class (higher threshold)
captures more served slots under FCFS.}{fig:balk-threshold}

\subsection{Demand Decomposition and Capacity Stress}

\mainfigure{fcfs_arrival_outcome_decomposition.png}{%
Full outcome decomposition stacked-bar view. Complements
Figure~\ref{fig:main-panel} with a more detailed rendering.
}{fig:decomp-full}

\mainfigure{fcfs_capacity_stress_curves.png}{%
Capacity stress diagnostic curves. Served rate and offered wait as
demand rises; utilization remains nearly flat.}{fig:stress}

\subsection{Balking Deep Dive: Additional Plots}

\mainfigure{../../../class1_balking/figures/mean_offered_delay_by_class.png}{%
Offered wait by class as Class~1 high-balking probability rises.
When enough Class~1 patients reject long-delay offers, overall horizon
congestion falls and both classes receive shorter offered delays.
}{fig:balk-offered}

\mainfigure{../../../class1_balking_threshold/figures/percent_serviced_by_class.png}{%
Class~1 balking threshold sweep: higher threshold increases Class~1
access by tolerating longer delays. The gain comes at mild cost to
Class~2 when Class~1 tolerance is very high.}{fig:thresh-served}

\mainfigure{../../../class1_balking_threshold/figures/service_gap.png}{%
Class~1 minus Class~2 service gap as Class~1 threshold varies.
}{fig:thresh-gap}

%% ============================================================
\section{Exact Metric Definitions}
\label{app:defs}
%% ============================================================

Let $i$ be patient class, $D$ measured service days, $S$ slots per day,
$A_i$ arrivals, $B_i$ accepted bookings, $K_i$ balked patients, $C_i$
cancellations, $H_i$ no-shows, $V_i$ completed visits, and $O_i = B_i +
K_i$ offered patients.

\begin{align*}
\texttt{percent\_serviced}_i &= V_i / A_i, \\
\texttt{overall\_percent\_serviced} &= \textstyle\sum_i V_i \;/\; \textstyle\sum_i A_i, \\
\texttt{mean\_offered\_booking\_delay} &= \textstyle\sum_i \tau^{\text{off}}_i \;/\; \textstyle\sum_i O_i, \\
\texttt{mean\_accepted\_booking\_delay} &= \textstyle\sum_i \tau^{\text{acc}}_i \;/\; \textstyle\sum_i B_i, \\
\texttt{average\_utilization} &= \tfrac{1}{D}\textstyle\sum_d V_d / S, \\
\texttt{balking\_rate} &= \textstyle\sum_i K_i \;/\; \textstyle\sum_i O_i, \\
\Delta_{\text{access}} &= \texttt{percent\_serviced}_1 - \texttt{percent\_serviced}_2.
\end{align*}

\noindent\textbf{Measurement semantics.} Class metrics cover patients
arriving in the measurement window. If \texttt{cooldown\_days}~$<$~%
\texttt{horizon\_days}~$-$~1, some late bookings are unresolved at
simulation end and counted as \texttt{unresolved\_booked} (0.06\% of
arrivals at baseline). \texttt{total\_value} is service-day based and
is not cohort-aligned with class metrics.

\end{document}
